% 论文基本信息
\ctitle{基于驱动程序替换的物联网设备固件仿真技术研究}
\etitle{Research on IoT Device Firmware Emulation Technology Based on Driver Replacement}

% 中文摘要
\cabstract{
随着物联网技术的快速发展，物联网设备数量呈现爆炸式增长，其固件安全问题日益突出。固件仿真技术是进行固件安全测试和漏洞挖掘的重要手段，然而现有的固件仿真方法存在外设建模成本高、硬件依赖强、自动化程度不足等问题。传统的基于硬件建模的方法需要为每种外设手动建立仿真模型，工作量大且难以覆盖多样化的硬件生态；基于驱动层替换的方法虽然避免了底层硬件建模，但驱动函数的定位与替换仍高度依赖人工干预，难以实现规模化应用。

针对上述问题，本文提出一种结合大语言模型的自动化固件仿真方法。该方法通过静态分析识别固件中的硬件依赖代码，利用大语言模型生成可替代的驱动实现，并通过动态测试反馈闭环迭代优化仿真质量，从而在无需手动建立外设模型的情况下支持固件稳定运行。围绕该方法，本文主要开展三方面工作。首先，设计并实现硬件依赖代码自动识别模块，从驱动程序源码中提取硬件相关函数调用点、外设寄存器访问模式和驱动函数边界，构建可查询的驱动知识库。其次，构建基于大语言模型的驱动函数分析与替换模块，结合结构化提示词模板和函数分类策略生成替代代码，并通过编译修复循环提升代码可编译性。最后，引入动态测试反馈闭环机制，在固件运行异常时自动分析原因并反馈至替换模块，通过"生成—测试—反馈—修正"流程持续优化仿真质量。

实验结果表明：在35个测试固件上，系统共识别4029个驱动函数并生成1563个替换实现。与现有代表性方法相比，本文方法的平均人工介入时长为0.6小时/固件，体现出较低的人力依赖和较高的自动化程度。上述结果说明本文方法可支持STM32、NXP、Atmel三类平台及FreeRTOS、RT-Thread、Zephyr、裸机等系统环境下的固件仿真，并在给定实验设置下具备较好的工程实用价值。
}

% 中文关键词
\ckeywords{物联网安全；固件仿真；驱动程序替换；大语言模型；静态分析}

% 英文摘要
\eabstract{
With the rapid development of Internet of Things (IoT) technologies, the number of IoT devices has grown explosively, and firmware security issues have become increasingly prominent. Firmware emulation is an important approach for firmware security testing and vulnerability discovery. However, existing firmware emulation methods still face major challenges, including high peripheral modeling cost, strong hardware dependence, and limited automation. Traditional hardware-modeling methods require manually constructing an emulation model for each peripheral, which is labor-intensive and difficult to scale to diverse hardware ecosystems. Although driver-layer replacement methods avoid low-level hardware modeling, driver-function identification and replacement still rely heavily on manual effort, limiting large-scale application.

To address these issues, this paper proposes an automated firmware emulation method enhanced by large language models (LLMs). The method identifies hardware-dependent firmware code through static analysis, generates alternative driver implementations with LLMs, and iteratively improves emulation quality through a dynamic testing feedback loop, thereby enabling stable firmware execution without manual peripheral modeling. Based on this method, this work makes three main contributions. First, it designs and implements an automatic hardware-dependency identification module that extracts hardware-related function call sites, peripheral register access patterns, and driver-function boundaries to build a queryable driver knowledge base. Second, it develops an LLM-based driver-function analysis and replacement module that combines structured prompts and function-classification strategies to generate replacement code and improve compilability through a compile-repair loop. Third, it introduces a dynamic testing feedback loop that analyzes runtime failures and feeds diagnostics back to the replacement module, continuously improving emulation quality through a "generate-test-feedback-revise" cycle.

Experimental results show that, across 35 firmware samples, the system identified 4,029 driver functions and generated 1,563 replacement implementations. Compared with representative baseline methods, the average manual intervention time was 0.6 hours per firmware, indicating lower human dependence and a higher degree of automation. These results indicate that the proposed method supports firmware emulation across STM32, NXP, and Atmel platforms, as well as FreeRTOS, RT-Thread, Zephyr, and bare-metal environments, and demonstrates practical engineering value under the given experimental setup.
}

% 英文关键词
\ekeywords{IoT Security; Firmware Emulation; Driver Replacement; Large Language Model; Static Analysis}
